\chapter{Marco Teórico}

% TODO: Párrafo introductorio que presente los conceptos fundamentales
% necesarios para comprender el resto del trabajo.

\section{Desarrollo Dirigido por Comportamiento}

% TODO: Definir BDD, origen, motivación. Explicar la sintaxis Gherkin
% (Given-When-Then). Rol en la ingeniería de software moderna.

\section{Pruebas Basadas en Modelos}

% TODO: Definir MBT. Explicar el proceso de generación automática de
% casos de prueba a partir de modelos formales.

\section{Sistemas de Transición Simbólicos}

% TODO: Definir STS formalmente. Explicar estados, transiciones,
% guardas, variables. Relación con MBT.

\section{Lenguaje Intermedio de BDD}

% TODO: Definir IBDD. Explicar la gramática formal (variables globales/locales,
% gates, guards, assignments). Rol como puente entre BDD y STS.
% Describir BDDTS y la cadena de transformación IBDD → BDDTS → STS.

\section{Modelos de Lenguaje Grandes}

% TODO: Describir LLMs relevantes para el trabajo. Capacidades de
% comprensión y generación de texto estructurado. Aplicaciones en
% ingeniería de software (generación de código, traducción de formalismos).

\section{Análisis Sintáctico y Gramáticas Formales}

% TODO: Conceptos de parsing necesarios para comprender el módulo de
% validación. Algoritmos relevantes (Earley, LALR). Notación EBNF.

\section{Trabajos Relacionados}

% TODO: Revisión de trabajos previos en: (1) automatización de BDD,
% (2) traducción de lenguaje natural a formalismos, (3) uso de LLMs
% para generación de especificaciones formales, (4) validación
% sintáctica de salidas de LLMs.
